\section{Improvements I: Data Integrity}
\label{sec:impdata}

Figure~\ref{fig:pipeline-new} shows the improved pipeline. Table~\ref{tab:changes} maps every
audit finding to its fix and the file that implements it. The next four sections describe the
changes in detail.

\begin{figure}[H]
\centering
\begin{tikzpicture}[node distance=0.35cm and 0.45cm,
  box/.style={draw,rounded corners,fill=white,align=center,font=\scriptsize,minimum height=0.95cm,text width=2.05cm},
  new/.style={box,fill=green!8,draw=green!45!black},
  arr/.style={-{Stealth},thick}]
\node[box] (csv) {\code{output.csv}};
\node[new,right=of csv] (clean) {parse numerics\\\textbf{deduplicate}};
\node[new,right=of clean] (bal) {balance\\(seeded rng)};
\node[box,right=of bal] (split) {80/20\\stratified split};
\node[new,right=of split] (prep) {train-only:\\log, z-score,\\PCA, $[0,\pi]$};
\node[new,below=0.9cm of prep] (zz) {ZZ / $R_x$ map\\bandwidth $\alpha$\\linear/circ./full};
\node[new,left=of zz] (ker) {fast exact\\statevectors\\$K=|\Phi\Phi^\dagger|^2$};
\node[new,left=of ker] (cv) {5-fold CV\\on $K_\text{train}$:\\$\alpha$, $C$};
\node[new,left=of cv] (eval) {metrics, McNemar,\\diagnostics,\\multi-seed};
\node[new,left=of eval] (rbf) {RBF-SVM\\CV: $C$, $\gamma$\\(+ RF ref.)};
\coordinate (L) at ([xshift=-4mm]csv.west);
\draw[arr] (csv)--(clean); \draw[arr] (clean)--(bal); \draw[arr] (bal)--(split); \draw[arr] (split)--(prep);
\draw[arr] (prep)--(zz); \draw[arr] (zz)--(ker); \draw[arr] (ker)--(cv); \draw[arr] (cv)--(eval);
\draw[arr] (rbf)--(eval);
\draw[arr,dashed] (prep.north) -- ++(0,0.35) node[above left,font=\scriptsize,pos=1]{} -| (L) |- (rbf.west);
\node[font=\scriptsize,above] at ($(prep.north)!0.5!(csv.north)+(0,0.35)$) {same PCA features};
\end{tikzpicture}
\caption{Improved pipeline. Green boxes are new or changed. Every fitted statistic comes from the
training split only.}
\label{fig:pipeline-new}
\end{figure}

\begin{table}[H]
\centering\small
\caption{Audit finding $\to$ change $\to$ location.}
\label{tab:changes}
\begin{tabularx}{\linewidth}{c X l}
\toprule
\# & Change & Where \\
\midrule
I1 & Drop exact duplicate (features, label) rows before sampling; report the train/test overlap & \code{data\_loader.py} \\
I2 & 5-fold CV grid search on the training set for QSVM ($\alpha$, $C$) and RBF ($C$, $\gamma$), same metric and folds & \code{evaluation.py} \\
I3 & Bandwidth $\alpha$ (input scaling) in all feature maps; concentration diagnostics & \code{qsvm.py}, \code{evaluation.py} \\
I4 & One simulation per sample plus a Gram product; fast exact NumPy engine; state cache & \code{qsvm.py} \\
I5 & All scaling/PCA fitted on training rows only; test values clipped to $[0,\pi]$ & \code{preprocessing.py} \\
I6 & Signed $\log(1+|x|)$ before standardisation and PCA & \code{preprocessing.py} \\
I7 & \code{--seeds}: repeat over seeds, report mean $\pm$ std; exact McNemar test per seed & \code{main.py}, \code{evaluation.py} \\
I8 & Seeded \code{default\_rng}; seeded Platt scaling, PCA and CV & all modules \\
I9 & All methods share one loader, one split and one preprocessing per seed & \code{main.py}, \code{benchmark.py} \\
I10 & \code{--kernel rx} (closed form) and \code{--sampling} (shots) added to the pipeline & \code{qsvm.py} \\
I11 & Balanced accuracy, macro-F1, MCC, AUC (binary and OvR), false-alarm and detection rates & \code{evaluation.py} \\
I12 & No padding; real feature counts reported; no silent truncation & \code{data\_loader.py}, \code{main.py} \\
I13 & Balance from the full (deduplicated) pool, so \code{--samples} is the real sample size & \code{data\_loader.py} \\
\bottomrule
\end{tabularx}
\end{table}

\subsection{Duplicate removal (I1)}
\paragraph{Problem.} With duplicates kept, many test flows are byte-identical to training flows.
\code{data\_loader.duplicate\_overlap()} measures this directly by hashing every training row and
counting test rows found in the set. For the original protocol it reports
\textbf{28--30\%} of test rows in the binary task (Table~\ref{tab:ablation}). A kernel
classifier gets those rows right automatically: $K(\bm x,\bm x)=1$ is the largest possible kernel
value, so the training twin dominates the decision function.

\paragraph{Change.} \code{load\_csv(dedup=True)} (the default) removes rows whose 61 feature values
\emph{and} task label equal those of an earlier row, before any sampling:
\begin{verbatim}
    key = X_df.copy(); key["__label"] = y_all
    keep = ~key.duplicated().values
\end{verbatim}
The key uses the label of the task (\code{y\_all}), not the raw attack name: in binary mode an
LDAP flow and an identical Portmap flow are both ``DDoS'', so the second one is a duplicate
(175 extra rows). Identical feature vectors with \emph{different} task labels (177 rows in
multiclass mode, 2 in binary mode) are kept. They are genuine
ambiguity, not leakage. \code{--keep-duplicates} restores the old behaviour for comparison.
After deduplication the overlap falls to 0--0.25\% (a handful of rows that differ only in
dropped zero-variance columns).

\paragraph{Why it matters for the conclusions.} Duplicates do not affect all models equally. A
kernel that is very sharp ($K\approx I$ off the diagonal, as with the untuned 16-qubit ZZ map)
behaves almost like a nearest-neighbour lookup. It benefits \emph{most} from exact twins and is
hurt most on genuinely new flows. Leakage can therefore create a spurious advantage for exactly
the most ``quantum'' kernel.

\subsection{Leakage-free preprocessing (I5)}
\paragraph{Problem.} The original loader called \code{MinMaxScaler().fit\_transform(X)} on all
sampled rows, then \code{main.py} split them. Each feature's minimum and maximum therefore
included test rows. With heavy-tailed features the maximum is often a single extreme flow, so
whether that flow lands in train or test changes the scale of every training feature.

\paragraph{Change.} \code{data\_loader.py} now returns raw, unscaled features. A new class,
\code{QuantumPreprocessor} in \code{preprocessing.py}, is fitted on the training split only, inside
\code{QSVM.fit()}. Its \code{transform()} applies the stored statistics to the test set and clips
out-of-range angles to $[0,\pi]$. The classical baselines use \code{transform\_pca()} of the
\emph{same} fitted object, so both models see identical information. The original behaviour is
kept as \code{LegacyPreprocessor} plus \code{legacy\_global\_scale()}, selected with
\code{--legacy-scaling}.

\subsection{Signed log transform (I6)}
\paragraph{Problem.} Angle encoding maps each feature to a rotation. If 90\% of flows fall within
0.002\,rad of each other, the corresponding qubit is in essentially the same state for all of
them, and the kernel cannot tell them apart along that direction. PCA does not help, because it
is dominated by the few extreme flows that define the variance.

\paragraph{Change.} Before standardisation and PCA, every feature is transformed by
\begin{equation}
x\mapsto \operatorname{sign}(x)\,\log(1+|x|).
\end{equation}
This is monotone, defined for zero and negative values (\code{Fwd Header Length} is negative in
12{,}592 rows, an artefact of CICFlowMeter), and compresses nine orders of magnitude into about 22 units.
Figure~\ref{fig:skew} shows the result for \code{Flow Bytes/s}: the median angle moves from
0.45 to 2.86\,rad and the fraction of flows below 0.01\,rad drops from 35\% to 4\%. The pipeline
is then: drop zero-variance columns $\to$ signed log $\to$ \code{StandardScaler} $\to$
\code{PCA}($n$) $\to$ per-component min--max to $[0,\pi]$, all fitted on training data.
\code{--no-log} disables the transform for ablation.

\subsection{Reproducible sampling and real sample sizes (I8, I12, I13)}
All randomness now flows from the \code{--seeds} values: \code{np.random.default\_rng(seed)} for
balancing, \code{random\_state=seed} for the split, PCA, \code{StratifiedKFold}, Platt scaling, the
Random Forest and the shot-noise generator. Two runs with the same seed produce byte-identical
results; this is verified in \code{tests/test\_core.py::test\_dedup\_and\_reproducible}. Classes are
balanced by drawing from the whole (deduplicated) pool of each class, so \code{--samples 2000} now
really gives 2000 samples (1000 per class in binary mode, 333 per class in multiclass mode)
instead of 822. \code{--legacy} uses the same sampler, so that the original and improved protocols
are always compared on the same number of flows. Numeric
parsing uses \code{pd.to\_numeric(errors='coerce')} so stray strings cannot crash the cast. Zero
padding to 58 columns was removed, and the true number of features entering PCA is printed.
